\documentclass{article}
\usepackage{graphicx} % Required for inserting images
\usepackage[margin=3cm]{geometry}
\usepackage{caption}
\usepackage{natbib}
\bibliographystyle{apalike}

\title{Banking on Biodiversity: Dietary Diversity Idea A Mockup}
\author{Carl Dinkel\thanks{Stefano Tarroni and "Banking on Biodiversity" Research team}\\University of Gothenburg}

\date{January 27, 2026}

\begin{document}

\maketitle

\begin{abstract}
    This document presents an exemplary research design intended as addition to the research project Banking on Biodiversity which focuses on the role of Community Seed Banks (CSBs) in shaping crop choices, resilience, and local comparative advantage. The core background, theoretical framework, and empirical methods are retained. This paper examines empirical evidence of nutritional aspects from various Development Economics literary sources and means to integrate the framework with that used for Banking on Biodiversity i.e., resilience mechanisms and supported agrobiodiversity that arises from expanded crop diversification, yield stability, adoption of improved seeds and strengthened household income to determine the effects of Community Seed Banks on household dietary diversity using a transformed Household Dietary Diversity Score (HDDS) style approach including deviations from the EAT-Lancet reference diet composition.
\end{abstract}

\section{Background and Motivation}

Community Seed Banks (CSBs) are farmer-led institutions that store, exchange and reproduce
locally adapted seeds, providing resilience against climate shocks and supporting agrobiodiversity.
In Uganda, CSBs have expanded over the last decade through programs supported by NARO,
Bioversity International, and several NGOs. Despite their growing presence, evidence remains
scarce on their economic and ecological impacts.

The Banking on Diversity project aims to quantify how access to CSBs affects:\begin{itemize}
    \item Crop and variety diversification;
    \item Productivity and yield stability;
    \item Adoption of improved and traditional seeds;
    \item Household income resilience and market orientation.    
\end{itemize}

The central hypothesis is that CSBs relax seed-access constraints that otherwise exclude farmers
from cultivating high-yield or stress-tolerant varieties, shaping local comparative advantage under
climate change.

The additional hypothesis that this paper is built upon is that the same relaxation of constraint resulting in improved income stability and seed-availability enabling consumption-smoothing more in line with the households intertemporal preferences also enables improved "nutrition-smoothing".

An important distinction is that households may smooth calories but not micronutrients, meaning the aspired quantification is the resiliency of dietary composition in households where CSBs are present as opposed to representative households where they are not. 

\section{Theory and Conceptual Framework}

\subsection{Literature}

The framework for the additional elements of the study focusing on nutritional resiliency are based on the following publications:\begin{itemize}
    \item \textbf{Strauss \& Thomas (1998)}, \textit{Health, Nutrition, and Economic Development}. Strauss and Thomas provide a foundational framework for understanding health and nutrition as endogenous components of economic behavior in developing economies. Rather than treating health as a purely biological or exogenous outcome, the authors emphasize its dual role as both a consumption good and a productive input, linking nutritional status directly to labor productivity, wages, and income. In this sense, nutrition operates as a feedback mechanism within the economy, closely aligned with efficiency wage models in which worker productivity depends on health investments. A central contribution of the paper is the recognition that credit and liquidity constraints, particularly severe in poor, rural, and agriculturally dependent settings can prevent households from making health and nutrition investments that would be privately and socially efficient. 
    Importantly, Strauss and Thomas highlight persistent gaps in the literature regarding the dynamic relationship between economic shocks, institutional environments, and health outcomes in low-income settings. They emphasize the lack of longitudinal evidence capturing how nutrition responds over time to income volatility, seasonality, and incomplete markets, and raise the possibility that inefficiencies arise in household health input allocations when access to credit and insurance is limited. These concerns are particularly salient in rural agrarian economies, where exposure to shocks and seasonal food scarcity is high.
    Building on these insights, this study examines community seed banks (CSBs) as locally embedded institutions that may relax the credit and insurance constraints emphasized by Strauss and Thomas. By stabilizing access to seeds across seasons and following adverse shocks, CSBs may facilitate smoother agricultural production and more stable dietary composition, affecting nutrition through channels beyond income alone. Studying the relationship between seed bank access and dietary outcomes therefore responds directly to calls in the literature to identify institutional mechanisms through which health investments become feasible in constrained environments, with implications for nutrition-sensitive agricultural policy and long-run welfare.
    \item \textbf{Headey et al. (2023)}, \textit{Measuring the affordability of healthy diets: Evidence from national household surveys.} Headey et al. develop a benchmarked approach to assessing diet quality in low-income food systems by comparing nationally representative household consumption patterns in Ethiopia, Kenya, Tanzania, and Uganda to the EAT-Lancet reference diet. Using survey-based food consumption estimates expressed in kilocalories per adult-equivalent per day, they decompose total intake into major food groups and quantify “gaps” relative to EAT-Lancet recommended intakes across locations and expenditure quintiles.

    \begin{center}
\includegraphics[width=0.75\linewidth]{Headey1.png}
\captionof{figure}{Food consumption calories (kcal/day) per adult-equivalent by major food groups in selected Sub-Saharan African countries relative to the EAT-Lancet reference diet. Source: Headey et al. (2023).}
\end{center}

\begin{center}
\includegraphics[width=0.75\linewidth]{Headey2.png}
\captionof{figure}{Percentage gaps between estimated household consumption and EAT-Lancet reference diet intakes by expenditure quintile in rural and urban Uganda. Source: Authors’ estimates reported in Headey et al. (2023).}
\end{center}

Two patterns are especially relevant for Uganda: (i) diets are heavily concentrated in starchy staples, and (ii) shortfalls are large for food groups typically associated with micronutrient adequacy and diet quality, such as fruits, vegetables, and pulses/nuts/seeds. Headey et al. argue that divergence from the reference diet reflects not only poverty and affordability constraints but also high relative prices and demand-side preference patterns for specific healthy foods. This benchmarking framework is useful for the present study because it motivates an outcome that captures diet composition—not only total calories—and provides a transparent reference point for interpreting whether institutional changes translate into nutritionally meaningful shifts in household diets.
        \item \textbf{Sahn (1989)}, \textit{Seasonal variability in Third World agriculture: The Consequences for Food Security}. 
    \item \textbf{Dercon \& Krishnan (2000)}, \textit{Vulnerability, seasonality and poverty in Ethiopia}.
    \item \textbf{Alderman, Hoddinott \& Kinsey (2006)}, \textit{Long term consequences of early childhood malnutrition}.
    \item \textbf{Devereux (1993)}, \textit{Theories of Famine}.    
\end{itemize}

\section{Research Questions}

Specific focus for this research design:\begin{itemize}
    \item \textbf{Welfare and food security:} Whats the effect of CSBs on agricultural income, yields, and household food security and diversity?
\end{itemize}

\section{Outcomes of Interest}

Specific outcome of interest for this research design:\begin{itemize}
    \item \textbf{Food security and dietary diversity:} Household food security index and dietary
diversity from consumption data.
\end{itemize}

\subsubsection{Outcome motivation and construction}

The following figures are made with values from Headey et al. (2023) and are used to visualize and contextualize the differential approach later discussed in Methods and Empirical Strategy section. 

\begin{center}
\includegraphics[width=0.95\linewidth]{EATLANCETTABLE.png}
\captionof{figure}{"Table I made while examining different ideas, sorry for the quality I made it on my Ipad while riding the bus, will fix a better one ASAP"}
\end{center}

Baseline deviations from the EAT-Lancet reference diet in Uganda highlight large excess consumption of starchy staples and deficits in pulses, fruits, and vegetables, motivating a distance-based measure of diet quality.
\begin{center}
\includegraphics[width=0.95\linewidth]{Differencesgraph.png}
\captionof{figure}{Graph made to illustrate potential differences in year t. Treatment and Control lines are entirely hypothesized to build context and intuition for outcome variable later discussed.}
\end{center}

\section{Methods and Empirical Strategy}

\subsection{Dietary Diversity Idea}

To identify the causal impact of CSBs on farmer dietary diversity, I will employ a difference-indifferences (DiD) approach.

Treatment is defined spatially: a household is considered exposed in year t if at least one CSB is active within a radius r kilometers of the household’s location in that year. Following the main project design, alternative radii are used as robustness checks, and specifications excluding households in buffer (“donut”) zones around treated areas are employed to mitigate spillover concerns. 

\subsubsection{Outcome variable: EAT-Lancet diet adherence}

The primary outcome is a continuous diet-quality index measuring household alignment with the EAT-Lancet reference diet. Using detailed household food consumption data from the Uganda National Panel Survey (UNPS), total household food intake is aggregated over a seven-day recall period and converted into adult-equivalent daily consumption. Individual food items are mapped into major EAT-Lancet food groups, including starchy staples; pulses, nuts, and seeds; fruits; vegetables; animal-source foods; dairy; added fats; and added sugars.

Diet quality is measured as the distance between household calorie shares and the EAT-Lancet recommended shares. Formally, for household $i$ in community $c$ at time $t$, the distance index is defined as:
\[
D_{ict} = \sum_{g} \left| s_{g,ict} - s_{g}^{EAT} \right|,
\]
where $s_{g,ict}$ denotes the share of total dietary energy derived from food group $g$, and $s_{g}^{EAT}$ is the corresponding EAT-Lancet reference share. Lower values of $D_{ict}$ indicate closer adherence to the reference diet. This index captures changes in dietary composition rather than total caloric intake and is interpreted as an aggregate measure of diet quality.

Secondary outcomes include food-group-specific calorie shares and a household dietary diversity indicator based on the number of food groups consumed over the recall period.\vspace{1\baselineskip}

\textbf{Alternative distance approach:} As a robustness measure, an alternative weighted distance index is also considered, allowing deviations in specific food groups to receive greater emphasis. The weighted distance index is defined as:
\[
D^{W}_{ict} = \sum_{g} w_g \left( s_{g,ict} - s_{g}^{EAT} \right)^2,
\]
where $w_g \geq 0$ denotes a pre-specified weight for food group $g$. Higher weights assign greater importance to deviations from the reference diet in selected food groups, such as fruits, vegetables, or pulses.

Weights are normalized to sum to one and are held constant across households and time. The weighted index penalizes large deviations more strongly than the unweighted measure and allows sensitivity analysis with respect to the nutritional relevance of different food groups. Results using the weighted index are reported as robustness checks to assess whether estimated treatment effects are driven by improvements in nutritionally salient components of the diet.

\subsubsection{Difference-in-differences specification}

For transparency and comparability, baseline estimates are obtained from a two-way fixed-effects DiD specification:
\[
D_{ict} = \alpha + \beta \, CSB_{ct}(r) + \gamma_c + \delta_t + X_{ict}'\theta + \varepsilon_{ict},
\]
where $CSB_{ct}(r)$ is an indicator for CSB exposure within radius $r$, $\gamma_c$ are community fixed effects, $\delta_t$ are survey-year fixed effects, and $X_{ict}$ is a vector of time-varying household controls.

\subsubsection{Staggered adoption and ATT estimation}

Because CSBs are introduced at different points in time and treatment effects may be heterogeneous, the primary analysis implements a staggered DiD estimator following Callaway and Sant’Anna (2021). Households are grouped by the year in which they are first exposed to a CSB, and group-time average treatment effects on the treated ($ATT_{g,t}$) are estimated by comparing treated households to not-yet-treated (and, where applicable, never-treated) households.

Dynamic treatment effects are examined using event-study style estimates of $ATT_{g,t}$, allowing dietary responses to evolve over time following CSB introduction. This is particularly relevant given that changes in agricultural practices and diet composition may occur gradually rather than immediately.

\subsubsection{Inference and robustness}

Standard errors are clustered at the administrative or community level consistent with the main project design to account for spatial correlation in treatment exposure. Robustness checks include alternative treatment radii, exclusion of buffer zones around treated areas, distance-based treatment intensity measures, and placebo tests using pre-treatment periods to assess the plausibility of the parallel trends assumption.

\subsection{Data}

The analysis will leverage multiple data sources:\begin{itemize}
    \item \textbf{Uganda National Panel Survey (UNPS / LSMS-ISA):} A panel of households surveyed 8 times (2005, 2009, 2010, 2011, 2013, 2015, 2018, 2019) by the Uganda Bureau of Statistics in collaboration with the World Bank. This dataset includes detailed information on household demographics, agricultural production, input use, shocks, and consumption. The panel structure allows for constructing a dataset suitable for difference-in-differences (DiD) analysis. This data is geocoded, enabling matching with CSB locations.
  \begin{figure}[h]
    \centering
    \includegraphics[width=0.75\linewidth]{Data.png}
  \end{figure}
  \item \textbf{CSB Registry:} A registry of all 27 CSBs active in Uganda, compiled through collaboration with NARO. This registry will include names, start of activity year, and locations. It will furthermore include measures of treatment strength (number of varieties, revolved seed quantities, number of members and customers) and other relevant characteristics. It will be constructed through a planned data collection including Community Seed Bank survey and digitization of CSB ledgers. \begin{figure}[h]
    \centering
    \includegraphics[width=0.75\linewidth]{data2.png}
  \end{figure}
  \item\textbf{Environmental Shocks:}  High resolution weather data (for example from the CHIRPS dataset) to construct measures of rainfall shocks (excessive and insufficient) at the community level.
  \item\textbf{Landsat 7 Imagery:} The Landsat program provides multispectral imagery at 30m spatial resolution, suitable for land use and vegetation monitoring. Landsat 7 offers a consistent archive covering entirely the 2005–2019. These data will be used to construct vegetation indices (NDVI, EVI), classifying cropland and forest cover.
  \item\textbf{ESA Global Plant Functional Type (PFT) Map Series:} The PFT global dataset includes 14 different plant functional types, each describing the percentage cover (0–100 percent) at a spatial resolution of 300 m at the equator. These layers importantly include forest cover, shrubs, natural grasses, and managed cropland.
  \item\textbf{EOD – eBird Observation Dataset:} eBird provides geocoded observational data of bird species sightings across Uganda. This data, while strongly dependent on observer effort and paths, can be used to construct a proxy for ecosystem biodiversity (Callaghan, 2019). 

\end{itemize}

\section{Early Results}

\begin{itemize}
    \item Preliminary descriptive evidence suggests that CSB villages exhibit higher on-farm varietal diversity and a greater share of improved seeds relative to control villages.
    \item Early analysis indicates a positive correlation between CSB presence and market participation for legumes and cereals.
    \item Yields appear more stable across seasons in CSB-supported communities, suggesting improved resilience to rainfall variation.
    \item Farmers report shorter seed search distances and reduced reliance on purchased inputs.
\end{itemize}

\section{Next Steps}

\begin{itemize}
    \item Finalize data linkage between the CSB registry and household survey.
    \item Pre-register hypotheses and main outcome variables.
    \item Implement main estimation and heterogeneity analyses.
    \item Draft the policy brief and prepare dissemination workshops with NARO and Bioversity International.
\end{itemize}

\nocite{*}
\bibliography{references}

\end{document}
